% The abstract of the notes.  Three parts, in this order: Översikt (what
% the deck does and where it goes), Lärandemål (the learning objectives),
% Förkunskaper (what the student should have met before).  A Läsning part
% is optional.  Section headings are \emph, not \section: the abstract is
% one block of prose.
%
% PROJECT: the headings below are Swedish.  Translate them, and the LO
% wording, for a deck in another language.

\emph{Översikt}
One paragraph: what question the deck opens with, what it takes the
student through, and where it stops.  Name topics and deck titles, never
weeks, lecture order or course events.

\emph{Lärandemål}
Efter denna modul ska studenten kunna:
% PROJECT: keep the course's own outcome codes here as a comment, so the
% mapping survives without leaking into the student's text.
% Lo04, Lo05 ...
\begin{restatable}{lo}{ModuleLOFirst}\label{ModuleLOFirst}%
  What the student can do, in plain words, without restating the
  mechanism.  The label is \texttt{ModuleLOAspect}: the module's name
  followed by \texttt{LO} and the aspect, so the notes can restate the
  objective where the prose meets it with \verb|\ModuleLOFirst|.
\end{restatable}
\begin{restatable}{lo}{ModuleLOSecond}\label{ModuleLOSecond}%
  A second objective.  One capability per environment; a sentence that
  needs an "och" for two unrelated abilities is two objectives.
\end{restatable}
\begin{restatable}{lo}{ModuleLOThird}\label{ModuleLOThird}%
  A third objective.
\end{restatable}

\emph{Förkunskaper}
What the student should have met first: "Studenten bör ha tagit del av
föreläsningen \emph{Deck Title}", never "ha sett", and never a week or a
lecture number.  Say what that deck supplies that this one needs.
